% GENERATED FILE. Edit canonical lessons, then run npm run generate:books.
% canonical-source-hash: fnv1a64:5b332288ea4f73a3
% canonical-lessons: HI-C141-makkhan, HI-C141-panir, HI-C141-sahad, HI-C141-mirc, HI-C141-ata

\chapter{Butter, Paneer, Honey, Chilli, Flour}
\label{ch:hi-butter-paneer-honey-chilli-flour}
\hlchaptermodality{141}

\begin{chapteropening}
\textbf{By the end of this chapter:} \emph{I can say butter, paneer, honey, a chilli and flour.}

Complete the last lesson of chapter 141: I can say butter, paneer, honey, a chilli and flour.
\end{chapteropening}

\section[\texorpdfstring{\hi{मक्खन} (makkhan)}{मक्खन (makkhan)}]{\hi{मक्खन} (makkhan) — butter}

\label{lesson:HI-C141-makkhan}



\begin{quote}
Before the new one: say the Hindi for a coconut, then the Hindi for yoghurt.
\end{quote}

\subsection*{You'll want to know: \hi{मक्खन}}
\textbf{\hi{मक्खन}} — \emph{makkhan} — \textquotedblleft{}butter\textquotedblright{}.

\begin{cousinweb}
\textbf{\hi{मक्खन}} (\emph{makkhan}) is a tadbhava, a word inherited from Sanskrit and reshaped by sound change. It comes from \emph{mrakṣaṇa}, \textquotedblleft{}rubbing in; ointment, oil\textquotedblright{} (CDIAL 10378). Butter was first \textquotedblleft{}the thing rubbed in\textquotedblright{}.
\end{cousinweb}

\begin{grammarlens}[title={What you've built}]
One more word for this chapter.
\end{grammarlens}

\subsection*{Your turn}
\begin{itemize}
\raggedright
  \item \emph{Say it:} \emph{makkhan}
  \item \emph{Say it:} \emph{makkhan}, once more
  \item \emph{Say it:} say \emph{makkhan}
  \item \emph{Recall:} say the Hindi for a coconut, then the Hindi for yoghurt, then say \emph{makkhan} again
\end{itemize}

\begin{tcolorbox}[breakable,colback=teal!4,colframe=teal!35!black,title={Before you move on}]
What does \hi{मक्खन} mean? (\textquotedblleft{}Butter\textquotedblright{}.) Say it once more.
\end{tcolorbox}
\section[\texorpdfstring{\hi{पनीर} (panīr)}{पनीर (panīr)}]{\hi{पनीर} (panīr) — paneer, cottage cheese}

\label{lesson:HI-C141-panir}



\begin{quote}
Before the new one: say the Hindi for yoghurt, then the Hindi for butter.
\end{quote}

\subsection*{You'll want to know: \hi{पनीर}}
\textbf{\hi{पनीर}} — \emph{panīr} — \textquotedblleft{}paneer, cottage cheese\textquotedblright{}.

\begin{cousinweb}
\textbf{\hi{पनीर}} (\emph{panīr}) is a loan, usually traced to Persian \emph{panīr}, \textquotedblleft{}cheese\textquotedblright{}.
\end{cousinweb}

\begin{grammarlens}[title={What you've built}]
One more word for this chapter.
\end{grammarlens}

\subsection*{Your turn}
\begin{itemize}
\raggedright
  \item \emph{Say it:} \emph{panīr}
  \item \emph{Say it:} \emph{panīr}, once more
  \item \emph{Say it:} say \emph{panīr}
  \item \emph{Recall:} say the Hindi for yoghurt, then the Hindi for butter, then say \emph{panīr} again
\end{itemize}

\begin{tcolorbox}[breakable,colback=teal!4,colframe=teal!35!black,title={Before you move on}]
What does \hi{पनीर} mean? (\textquotedblleft{}Paneer, cottage cheese\textquotedblright{}.) Say it once more.
\end{tcolorbox}
\section[\texorpdfstring{\hi{शहद} (śahad)}{शहद (śahad)}]{\hi{शहद} (śahad) — honey}

\label{lesson:HI-C141-sahad}



\begin{quote}
Before the new one: say the Hindi for butter, then the Hindi for paneer.
\end{quote}

\subsection*{You'll want to know: \hi{शहद}}
\textbf{\hi{शहद}} — \emph{śahad} — \textquotedblleft{}honey\textquotedblright{}.

\begin{cousinweb}
\textbf{\hi{शहद}} (\emph{śahad}) is a loan from Arabic \emph{shahd}, \textquotedblleft{}honey\textquotedblright{}, which reached Hindi through Persian. Its Arabic root is \emph{sh-h-d} (Lane).
\end{cousinweb}

\begin{grammarlens}[title={What you've built}]
One more word for this chapter.
\end{grammarlens}

\subsection*{Your turn}
\begin{itemize}
\raggedright
  \item \emph{Say it:} \emph{śahad}
  \item \emph{Say it:} \emph{śahad}, once more
  \item \emph{Say it:} say \emph{śahad}
  \item \emph{Recall:} say the Hindi for butter, then the Hindi for paneer, then say \emph{śahad} again
\end{itemize}

\begin{tcolorbox}[breakable,colback=teal!4,colframe=teal!35!black,title={Before you move on}]
What does \hi{शहद} mean? (\textquotedblleft{}Honey\textquotedblright{}.) Say it once more.
\end{tcolorbox}
\section[\texorpdfstring{\hi{मिर्च} (mirc)}{मिर्च (mirc)}]{\hi{मिर्च} (mirc) — a chilli}

\label{lesson:HI-C141-mirc}



\begin{quote}
Before the new one: say the Hindi for paneer, then the Hindi for honey.
\end{quote}

\subsection*{You'll want to know: \hi{मिर्च}}
\textbf{\hi{मिर्च}} — \emph{mirc} — \textquotedblleft{}a chilli\textquotedblright{}.

\begin{cousinweb}
\textbf{\hi{मिर्च}} (\emph{mirc}) is a tadbhava, a word inherited from Sanskrit and reshaped by sound change. It comes from \emph{marīca}, \textquotedblleft{}peppercorn\textquotedblright{} (CDIAL 9875). CDIAL takes the Sanskrit word from an Austroasiatic language. When the chilli came to India, it took the name of black pepper.
\end{cousinweb}

\begin{grammarlens}[title={What you've built}]
One more word for this chapter.
\end{grammarlens}

\subsection*{Your turn}
\begin{itemize}
\raggedright
  \item \emph{Say it:} \emph{mirc}
  \item \emph{Say it:} \emph{mirc}, once more
  \item \emph{Say it:} say \emph{mirc}
  \item \emph{Recall:} say the Hindi for paneer, then the Hindi for honey, then say \emph{mirc} again
\end{itemize}

\begin{tcolorbox}[breakable,colback=teal!4,colframe=teal!35!black,title={Before you move on}]
What does \hi{मिर्च} mean? (\textquotedblleft{}A chilli\textquotedblright{}.) Say it once more.
\end{tcolorbox}
\section[\texorpdfstring{\hi{आटा} (āṭā)}{आटा (āṭā)}]{\hi{आटा} (āṭā) — flour}

\label{lesson:HI-C141-ata}



\begin{quote}
Before the new one: say the Hindi for honey, then the Hindi for a chilli.
\end{quote}

\subsection*{You'll want to know: \hi{आटा}}
\textbf{\hi{आटा}} — \emph{āṭā} — \textquotedblleft{}flour\textquotedblright{}.

\begin{cousinweb}
\textbf{\hi{आटा}} (\emph{āṭā}) is a tadbhava, a word inherited and reshaped by sound change. It comes from \emph{*ārta}, \textquotedblleft{}flour\textquotedblright{}, a form CDIAL reconstructs from its descendants (CDIAL 1338). CDIAL compares Persian \emph{ārd}, \textquotedblleft{}flour\textquotedblright{}.
\end{cousinweb}

\begin{grammarlens}[title={What you've built}]
One more word for this chapter.
\end{grammarlens}

\subsection*{Your turn}
\begin{itemize}
\raggedright
  \item \emph{Say it:} \emph{āṭā}
  \item \emph{Say it:} \emph{āṭā}, once more
  \item \emph{Say it:} say \emph{āṭā}
  \item \emph{Recall:} say the Hindi for honey, then the Hindi for a chilli, then say \emph{āṭā} again
\end{itemize}

\begin{tcolorbox}[breakable,colback=teal!4,colframe=teal!35!black,title={Before you move on}]
What does \hi{आटा} mean? (\textquotedblleft{}Flour\textquotedblright{}.) Say it once more.
\end{tcolorbox}
